\begin{center}
  {\Huge \color{mainblue}\textbf{Esercitazione: la somma di frazioni}}\\[0.2cm]
  \rule{0.82\linewidth}{0.5pt}
\end{center}


\section*{Richiamo iniziale}

\begin{esercizio}{Il cocktail di frutta}
Per festeggiare la fine dell'anno scolastico, la classe decide di preparare un cocktail di frutta analcolico. La ricetta prevede di mescolare i seguenti ingredienti:
\begin{itemize}
  \item $\frac{1}{2}\unita{l}$ di succo di mela;
  \item $\frac{1}{3}\unita{l}$ di succo d'arancia;
  \item $\frac{1}{6}\unita{l}$ di sciroppo di fragola.
\end{itemize}

\domanda{a) Rappresenta graficamente gli ingredienti nel cilindro graduato sottostante.}\\
\domanda{b) Calcola la frazione totale di liquido ottenuta.}\\
\domanda{c) Se nella classe ci sono 24 studenti e ciascuno beve $\frac{1}{4}\unita{l}$ di cocktail, quanti litri servono in totale? Quante volte bisogna preparare la ricetta base per soddisfare tutta la classe?}

\vspace{0.4cm}

\begin{center}
\begin{tikzpicture}[scale=1.1, every node/.style={scale=0.9}]
  % Cylinder outline
  \draw[thick, linegray] (0,0) -- (0,5) (3,0) -- (3,5);
  \draw[thick, linegray] (0,0) to[out=-20,in=200] (3,0);
  \draw[thick, linegray] (1.5,5) ellipse (1.5 and 0.25);
  
  % Graduations on the left side
  \foreach \x in {1,2,3,4,5,6} {
    \draw[linegray, thick] (0, 0.75*\x) -- (0.4, 0.75*\x);
    \node[anchor=east, textgray] at (-0.15, 0.75*\x) {\small $\frac{\x}{6}\unita{l}$};
  }
  
  % Filling with colors (conditional on solutions)
  \ifsoluzioni
    % 1/6 strawberry syrup (red) from y=0 to y=0.75
    \path[fill=red!20!white] (0,0) to[out=-20,in=200] (3,0) -- (3, 0.75) to[out=200,in=-20] (0, 0.75) -- cycle;
    \draw[<-] (1.5, 0.3) -- (3.8, 0.3) node[right, scale=0.8] {Sciroppo di fragola ($\frac{1}{6}\unita{l}$)};
    
    % 2/6 orange juice (orange) from y=0.75 to y=2.25
    \path[fill=orange!35!white] (0,0.75) to[out=-20,in=200] (3,0.75) -- (3, 2.25) to[out=200,in=-20] (0, 2.25) -- cycle;
    \draw[black] (1.5, 1.5) node {\small Succo d'arancia ($\frac{1}{3}\unita{l}$)};
    
    % 3/6 apple juice (yellow) from y=2.25 to y=4.5
    \path[fill=yellow!30!white] (0,2.25) to[out=-20,in=200] (3,2.25) -- (3, 4.5) to[out=200,in=-20] (0, 4.5) -- cycle;
    \draw[black] (1.5, 3.375) node {\small Succo di mela ($\frac{1}{2}\unita{l}$)};
    
    % Ellipses for liquid layers
    \draw[red!80!black, fill=red!30!white, opacity=0.8] (1.5,0.75) ellipse (1.5 and 0.15);
    \draw[orange!85!black, fill=orange!40!white, opacity=0.8] (1.5,2.25) ellipse (1.5 and 0.15);
    \draw[yellow!80!black!80!red, fill=yellow!30!white, opacity=0.8] (1.5,4.5) ellipse (1.5 and 0.15);
  \else
    % Dashed ellipses inside cylinder for students to fill in
    \foreach \x in {1,2,3,4,5} {
      \draw[linegray, dashed] (1.5, 0.75*\x) ellipse (1.5 and 0.15);
    }
  \fi
\end{tikzpicture}
\end{center}
 
\responsespace{5.5cm}
 
\begin{soluzione}
  Passaggi risolutivi:
  \begin{enumerate}[label=\alph*)]
    \item Rappresentazione grafica (vedi cilindro graduato suddiviso in sesti).
    \item Frazione totale di liquido: Sommiamo le quantità:
      \[
      \frac{1}{2} + \frac{1}{3} + \frac{1}{6} = \frac{3}{6} + \frac{2}{6} + \frac{1}{6} = \frac{3+2+1}{6} = \frac{6}{6} = 1\unita{l}
      \]
    \item Fabbisogno per la classe:
      \begin{itemize}
        \item Quantità totale per 24 studenti: $24 \cdot \frac{1}{4}\unita{l} = \frac{24}{4}\unita{l} = 6\unita{l}$.
        \item Poiché una singola ricetta base produce esattamente $1\unita{l}$ di liquido, per preparare $6\unita{l}$ totali sarà necessario ripetere la ricetta esattamente $6$ volte.
      \end{itemize}
  \end{enumerate}
\end{soluzione}
\end{esercizio}

\condnewpage

\begin{tcolorbox}[
  colback=blue!3,
  colframe=mainblue,
  title={Somma di frazioni},
  fonttitle=\bfseries\large,
  arc=2mm,
  boxrule=1pt,
  left=4mm, right=4mm, top=3mm, bottom=3mm
]
Per sommare due frazioni dobbiamo avere parti della stessa grandezza. \textbf{Perché dobbiamo trovare lo stesso denominatore?} Non possiamo sommare parti di misure diverse (come metà e terzi). Calcolando il \textbf{m.c.m.} (minimo comune multiplo) dividiamo gli interi in parti della \textbf{stessa identica misura}. Solo a questo punto possiamo sommare i numeratori!

\begin{center}
\begin{tikzpicture}[scale=0.95, every node/.style={scale=0.95}]
  % Bar 1 (1/2 = 3/6)
  \node[anchor=east] at (-0.2, 3.25) {\large $\frac{1}{2} = \frac{3}{6}$};
  \fill[blue!15!white] (0,3.0) rectangle (3,3.5);
  \draw[thick, mainblue] (0,3.0) rectangle (6,3.5);
  \draw[thick, mainblue] (3,3.0) -- (3,3.5);
  \draw[dashed, linegray] (1,3.0) -- (1,3.5) (2,3.0) -- (2,3.5) (4,3.0) -- (4,3.5) (5,3.0) -- (5,3.5);

  % Bar 2 (1/3 = 2/6)
  \node[anchor=east] at (-0.2, 1.75) {\large $\frac{1}{3} = \frac{2}{6}$};
  \fill[blue!15!white] (0,1.5) rectangle (2,2.0);
  \draw[thick, mainblue] (0,1.5) rectangle (6,2.0);
  \draw[thick, mainblue] (2,1.5) -- (2,2.0) (4,1.5) -- (4,2.0);
  \draw[dashed, linegray] (1,1.5) -- (1,2.0) (3,1.5) -- (3,2.0) (5,1.5) -- (5,2.0);

  % Bar 3 (3/6 + 2/6 = 5/6)
  \node[anchor=east] at (-0.2, 0.25) {\large $\frac{3}{6} + \frac{2}{6} = \frac{5}{6}$};
  \fill[blue!35!white] (0,0) rectangle (5,0.5);
  \draw[thick, mainblue] (0,0) rectangle (6,0.5);
  \foreach \x in {1,2,3,4,5} {
    \draw[thick, mainblue] (\x,0) -- (\x,0.5);
  }
\end{tikzpicture}
\end{center}

\vspace{0.15cm}\hrule\vspace{0.15cm}

Esempio di calcolo:
$$\frac{1}{2} + \frac{1}{3}$$
\begin{enumerate}[label=\arabic*., leftmargin=1.5cm, noitemsep]
  \item Denominatore comune (m.c.m. tra 2 e 3): \textbf{6}
  \item Adatta i numeratori:
    $$\frac{1}{2} = \frac{1 \cdot 3}{6} = \frac{3}{6} \qquad \text{e} \qquad \frac{1}{3} = \frac{1 \cdot 2}{6} = \frac{2}{6}$$
  \item Somma i numeratori:
    $$\frac{3}{6} + \frac{2}{6} = \frac{5}{6}$$
\end{enumerate}
\end{tcolorbox}

\condnewpage


\section*{1. Esercizi di riscaldamento}

\begin{esercizio}{La piramide delle somme}
In questa piramide, il valore di ogni mattone è uguale alla somma dei valori dei due mattoni che si trovano direttamente sotto di esso.

\domanda{Completa la piramide inserendo le frazioni mancanti nei mattoni vuoti. Semplifica sempre le frazioni ottenute.}

\vspace{0.4cm}
\begin{center}
\begin{tikzpicture}[scale=1.0]
  % Draw bottom row
  \draw[thick, mainblue] (0,0) rectangle (2.6,1.3);
  \draw[thick, mainblue] (2.6,0) rectangle (5.2,1.3);
  \draw[thick, mainblue] (5.2,0) rectangle (7.8,1.3);
  
  % Draw middle row
  \draw[thick, mainblue] (1.3,1.3) rectangle (3.9,2.6);
  \draw[thick, mainblue] (3.9,1.3) rectangle (6.5,2.6);
  
  % Draw top row
  \draw[thick, mainblue] (2.6,2.6) rectangle (5.2,3.9);
  
  % Labels for bottom row (always visible)
  \node at (1.3,0.65) {\large $\frac{1}{3}$};
  \node at (3.9,0.65) {\large $\frac{1}{6}$};
  \node at (6.5,0.65) {\large $\frac{2}{5}$};
  
  % Labels for middle and top rows
  \ifsoluzioni
    \node[red!80!black] at (2.6,1.95) {\large \textbf{$\frac{1}{2}$}};
    \node[red!80!black] at (5.2,1.95) {\large \textbf{$\frac{17}{30}$}};
    \node[red!80!black] at (3.9,3.25) {\large \textbf{$\frac{16}{15}$}};
  \else
    \node[lightgray] at (2.6,1.95) {\textbf{?}};
    \node[lightgray] at (5.2,1.95) {\textbf{?}};
    \node[lightgray] at (3.9,3.25) {\textbf{?}};
  \fi
\end{tikzpicture}
\end{center}

\responsespace{4.0cm}

\begin{soluzione}
  \textbf{Passaggi:}
  \begin{enumerate}
    \item \textbf{Mattone intermedio sinistro:} somma dei due mattoni sottostanti (sinistro e centrale):
      \[
      \frac{1}{3} + \frac{1}{6} = \frac{2}{6} + \frac{1}{6} = \frac{3}{6} = \frac{1}{2}
      \]
    \item \textbf{Mattone intermedio destro:} somma dei due mattoni sottostanti (centrale e destro):
      \[
      \frac{1}{6} + \frac{2}{5} = \frac{5}{30} + \frac{12}{30} = \frac{17}{30}
      \]
    \item \textbf{Mattone superiore (cima):} somma dei due mattoni intermedi:
      \[
      \frac{1}{2} + \frac{17}{30} = \frac{15}{30} + \frac{17}{30} = \frac{32}{30} = \frac{16}{15}
      \]
  \end{enumerate}
\end{soluzione}
\end{esercizio}

\condnewpage


\begin{esercizio}{Uguaglianze da completare}
Risolvi le seguenti uguaglianze determinando la frazione mancante da inserire nel quadretto grigio. Semplifica i risultati quando possibile.

\vspace{0.3cm}
\renewcommand{\arraystretch}{2.25}
\begin{tabularx}{\linewidth}{|X|X|}
\hline
a) $\frac{2}{3} + \tcbox[colback=gray!15,colframe=gray!60,on line,boxsep=2pt,arc=1pt]{\makebox[0.8cm]{\ifsoluzioni \color{red!80!black}$\mathbf{\frac{5}{3}}$ \else \phantom{$\frac{1}{2}$} \fi}} = \frac{7}{3}$ & 
b) $\tcbox[colback=gray!15,colframe=gray!60,on line,boxsep=2pt,arc=1pt]{\makebox[0.8cm]{\ifsoluzioni \color{red!80!black}$\mathbf{\frac{4}{15}}$ \else \phantom{$\frac{1}{2}$} \fi}} + \frac{3}{5} = \frac{13}{15}$ \\
\hline
c) $\tcbox[colback=gray!15,colframe=gray!60,on line,boxsep=2pt,arc=1pt]{\makebox[0.8cm]{\ifsoluzioni \color{red!80!black}$\mathbf{\frac{2}{5}}$ \else \phantom{$\frac{1}{2}$} \fi}} + \frac{1}{5} = \frac{3}{5}$ & 
d) $\frac{3}{14} + \tcbox[colback=gray!15,colframe=gray!60,on line,boxsep=2pt,arc=1pt]{\makebox[0.8cm]{\ifsoluzioni \color{red!80!black}$\mathbf{\frac{5}{14}}$ \else \phantom{$\frac{1}{2}$} \fi}} = \frac{4}{7}$ \\
\hline
\end{tabularx}
\vspace{0.3cm}

\responsespace{3.8cm}

\begin{soluzione}
  Ragionamento per completare i quadretti:
  \begin{enumerate}[label=\alph*)]
    \item \textbf{Stesso denominatore:}
      Ci chiediamo: quale frazione dobbiamo aggiungere a $\frac{2}{3}$ per ottenere $\frac{7}{3}$?
      Poiché le frazioni hanno lo stesso denominatore (3), il denominatore rimarrà 3. Dobbiamo solo trovare il numeratore: cosa dobbiamo aggiungere a 2 per ottenere 7? Dobbiamo aggiungere 5 (infatti $2 + 5 = 7$).
      Quindi, la frazione da inserire è $\frac{5}{3}$.
    
    \item \textbf{Denominatore diverso:}
      Ci chiediamo: quale frazione dobbiamo aggiungere a $\frac{3}{5}$ per ottenere $\frac{13}{15}$?
      Per ragionare facilmente, trasformiamo $\frac{3}{5}$ in una frazione equivalente con denominatore 15. Moltiplicando numeratore e denominatore per 3, otteniamo $\frac{9}{15}$.
      Ora ci chiediamo: cosa dobbiamo aggiungere a $\frac{9}{15}$ per ottenere $\frac{13}{15}$?
      Poiché ora hanno lo stesso denominatore, dobbiamo semplicemente aggiungere 4 al numeratore (infatti $9 + 4 = 13$).
      Quindi, la frazione da inserire è $\frac{4}{15}$.
      
    \item \textbf{Stesso denominatore:}
      Ci chiediamo: quale frazione dobbiamo aggiungere a $\frac{1}{5}$ per ottenere $\frac{3}{5}$?
      Poiché hanno lo stesso denominatore (5), dobbiamo aggiungere 2 al numeratore (infatti $2 + 1 = 3$).
      Quindi, la frazione da inserire è $\frac{2}{5}$.
      
    \item \textbf{Denominatore diverso:}
      Ci chiediamo: quale frazione dobbiamo aggiungere a $\frac{3}{14}$ per ottenere $\frac{4}{7}$?
      Trasformiamo $\frac{4}{7}$ in una frazione equivalente con denominatore 14. Moltiplicando numeratore e denominatore per 2, otteniamo $\frac{8}{14}$.
      Ora ci chiediamo: cosa dobbiamo aggiungere a $\frac{3}{14}$ per ottenere $\frac{8}{14}$?
      Dobbiamo aggiungere 5 al numeratore (infatti $3 + 5 = 8$).
      Quindi, la frazione da inserire è $\frac{5}{14}$.
  \end{enumerate}
\end{soluzione}
\end{esercizio}

\begin{esercizio}{Corrispondenze tra operazioni e risultati}
Associa ciascuna operazione a sinistra con il suo risultato corretto a destra collegandoli con una freccia.

\vspace{0.4cm}
\begin{minipage}{0.48\linewidth}
  \begin{enumerate}[label=\arabic*)]
    \item $\frac{1}{2} + \frac{1}{4}$
    \vspace{0.4cm}
    \item $\frac{1}{3} + \frac{1}{6}$
    \vspace{0.4cm}
    \item $\frac{3}{5} + \frac{1}{10}$
    \vspace{0.4cm}
    \item $\frac{1}{12} + \frac{1}{4}$
  \end{enumerate}
\end{minipage}
\hfill
\begin{minipage}{0.48\linewidth}
  \begin{enumerate}[label=\Alph*)]
    \item $\frac{1}{3}$
    \vspace{0.4cm}
    \item $\frac{3}{4}$
    \vspace{0.4cm}
    \item $\frac{1}{2}$
    \vspace{0.4cm}
    \item $\frac{7}{10}$
  \end{enumerate}
\end{minipage}
\vspace{0.4cm}

\responsespace{2.2cm}

\begin{soluzione}
  Associazioni corrette:
  \begin{itemize}
    \item 1) $\frac{1}{2} + \frac{1}{4} = \frac{2}{4} + \frac{1}{4} = \frac{3}{4}$ $\to$ si associa alla lettera B)
    \item 2) $\frac{1}{3} + \frac{1}{6} = \frac{2}{6} + \frac{1}{6} = \frac{3}{6} = \frac{1}{2}$ $\to$ si associa alla lettera C)
    \item 3) $\frac{3}{5} + \frac{1}{10} = \frac{6}{10} + \frac{1}{10} = \frac{7}{10}$ $\to$ si associa alla lettera D)
    \item 4) $\frac{1}{12} + \frac{1}{4} = \frac{1}{12} + \frac{3}{12} = \frac{4}{12} = \frac{1}{3}$ $\to$ si associa alla lettera A)
  \end{itemize}
\end{soluzione}
\end{esercizio}

\condnewpage

\begin{tcolorbox}[
  colback=blue!3,
  colframe=mainblue,
  title={Sottrazione di frazioni},
  fonttitle=\bfseries\large,
  arc=2mm,
  boxrule=1pt,
  left=4mm, right=4mm, top=3mm, bottom=3mm
]
Anche per sottrarre due frazioni dobbiamo avere parti della stessa grandezza. \textbf{Perché dobbiamo trovare lo stesso denominatore?} Non possiamo sottrarre parti di misure diverse (come terzi e sesti). Calcolando il \textbf{m.c.m.} (minimo comune multiplo) dividiamo gli interi in parti della \textbf{stessa identica misura}. Solo a questo punto possiamo sottrarre i numeratori!

\begin{center}
\begin{tikzpicture}[scale=0.95, every node/.style={scale=0.95}]
  % Bar 1 (2/3 = 4/6)
  \node[anchor=east] at (-0.2, 3.25) {\large $\frac{2}{3} = \frac{4}{6}$};
  \fill[blue!15!white] (0,3.0) rectangle (4,3.5);
  \draw[thick, mainblue] (0,3.0) rectangle (6,3.5);
  \draw[thick, mainblue] (2,3.0) -- (2,3.5) (4,3.0) -- (4,3.5);
  \draw[dashed, linegray] (1,3.0) -- (1,3.5) (3,3.0) -- (3,3.5) (5,3.0) -- (5,3.5);

  % Bar 2 (1/6)
  \node[anchor=east] at (-0.2, 1.75) {\large $\frac{1}{6}$};
  \fill[blue!15!white] (0,1.5) rectangle (1,2.0);
  \draw[thick, mainblue] (0,1.5) rectangle (6,2.0);
  \foreach \x in {1,2,3,4,5} {
    \draw[thick, mainblue] (\x,1.5) -- (\x,2.0);
  }

  % Bar 3 (4/6 - 1/6 = 3/6 = 1/2)
  \node[anchor=east] at (-0.2, 0.25) {\large $\frac{4}{6} - \frac{1}{6} = \frac{3}{6} = \frac{1}{2}$};
  \fill[blue!35!white] (0,0) rectangle (3,0.5);
  \draw[thick, mainblue] (0,0) rectangle (6,0.5);
  \foreach \x in {1,2,3,4,5} {
    \draw[thick, mainblue] (\x,0) -- (\x,0.5);
  }
\end{tikzpicture}
\end{center}

\vspace{0.15cm}\hrule\vspace{0.15cm}

Esempio di calcolo:
$$\frac{2}{3} - \frac{1}{6}$$
\begin{enumerate}[label=\arabic*., leftmargin=1.5cm, noitemsep]
  \item Denominatore comune (m.c.m. tra 3 e 6): \textbf{6}
  \item Adatta i numeratori:
    $$\frac{2}{3} = \frac{2 \cdot 2}{6} = \frac{4}{6}$$
  \item Sottrai i numeratori:
    $$\frac{4}{6} - \frac{1}{6} = \frac{3}{6} = \frac{1}{2}$$
\end{enumerate}
\end{tcolorbox}

\condnewpage


\section*{2. Palestra di calcolo}

\begin{esercizio}{Calcoli con frazioni}
Calcola il risultato delle seguenti somme e sottrazioni tra frazioni, riducendo sempre i risultati ai minimi termini.

\vspace{0.2cm}
\renewcommand{\arraystretch}{2.5}
\begin{tabularx}{\linewidth}{XX}
  a) $\frac{5}{12} + \frac{1}{12} = \fillbox{2.2cm}$ & b) $\frac{7}{10} - \frac{3}{10} = \fillbox{2.2cm}$ \\
  c) $\frac{3}{4} + \frac{1}{6} = \fillbox{2.2cm}$ & d) $\frac{5}{6} - \frac{3}{8} = \fillbox{2.2cm}$ \\
  e) $1 - \frac{3}{7} = \fillbox{2.2cm}$ & f) $\frac{9}{4} - 2 = \fillbox{2.2cm}$ \\
  g) $3 - \frac{5}{6} = \fillbox{2.2cm}$ & h) $\frac{11}{3} - 3 = \fillbox{2.2cm}$ \\
  i) $\frac{24}{25} - \frac{11}{20} - \frac{1}{4} = \fillbox{2.2cm}$ & j) $\frac{1}{5} + \frac{21}{30} - \frac{1}{6} = \fillbox{2.2cm}$ \\
  k) $5 - \frac{1}{2} + \frac{1}{3} - \frac{1}{6} + \frac{1}{7} = \fillbox{2.2cm}$ & l) $\frac{44}{45} + \frac{17}{30} - \frac{2}{9} + \frac{1}{15} = \fillbox{2.2cm}$ \\
  m) $\frac{35}{50} - \frac{15}{16} + \frac{3}{40} - \frac{5}{200} = \fillbox{2.2cm}$ & \\
\end{tabularx}
\vspace{0.2cm}

\responsespace{3.2cm}

\begin{soluzione}
  Svolgimento passaggi:
  \begin{enumerate}[label=\alph*)]
    \item Stesso denominatore: $\frac{5}{12} + \frac{1}{12} = \frac{5+1}{12} = \frac{6}{12} = \frac{1}{2}$
    \item Stesso denominatore: $\frac{7}{10} - \frac{3}{10} = \frac{7-3}{10} = \frac{4}{10} = \frac{2}{5}$
    \item Comune denominatore (12): $\frac{3}{4} + \frac{1}{6} = \frac{9}{12} + \frac{2}{12} = \frac{11}{12}$
    \item Comune denominatore (24): $\frac{5}{6} - \frac{3}{8} = \frac{20}{24} - \frac{9}{24} = \frac{11}{24}$
    \item Intero e frazione: $1 - \frac{3}{7} = \frac{7}{7} - \frac{3}{7} = \frac{4}{7}$
    \item Frazione e intero: $\frac{9}{4} - 2 = \frac{9}{4} - \frac{8}{4} = \frac{1}{4}$
    \item Intero e frazione: $3 - \frac{5}{6} = \frac{18}{6} - \frac{5}{6} = \frac{13}{6}$
    \item Frazione e intero: $\frac{11}{3} - 3 = \frac{11}{3} - \frac{9}{3} = \frac{2}{3}$
    \item Comune denominatore (100):
      \[
      \frac{24}{25} - \frac{11}{20} - \frac{1}{4} = \frac{24 \cdot 4 - 11 \cdot 5 - 1 \cdot 25}{100} = \frac{96 - 55 - 25}{100} = \frac{16}{100} = \frac{4}{25}
      \]
    \item Semplificazione e mcm (30): Semplifichiamo prima $\frac{21}{30} = \frac{7}{10}$:
      \[
      \frac{1}{5} + \frac{7}{10} - \frac{1}{6} = \frac{1 \cdot 6 + 7 \cdot 3 - 1 \cdot 5}{30} = \frac{6 + 21 - 5}{30} = \frac{22}{30} = \frac{11}{15}
      \]
    \item Comune denominatore (42):
      \[
      5 - \frac{1}{2} + \frac{1}{3} - \frac{1}{6} + \frac{1}{7} = \frac{210 - 21 + 14 - 7 + 6}{42} = \frac{202}{42} = \frac{101}{21}
      \]
    \item Comune denominatore (90):
      \[
      \frac{44}{45} + \frac{17}{30} - \frac{2}{9} + \frac{1}{15} = \frac{88 + 51 - 20 + 6}{90} = \frac{125}{90} = \frac{25}{18}
      \]
    \item Semplificazione e mcm (80): Semplifichiamo $\frac{35}{50} = \frac{7}{10}$ e $\frac{5}{200} = \frac{1}{40}$:
      \[
      \frac{7}{10} - \frac{15}{16} + \frac{3}{40} - \frac{1}{40} = \frac{7}{10} - \frac{15}{16} + \frac{2}{40} = \frac{7}{10} - \frac{15}{16} + \frac{1}{20}
      \]
      Calcoliamo con denominatore comune 80:
      \[
      \frac{7 \cdot 8 - 15 \cdot 5 + 1 \cdot 4}{80} = \frac{56 - 75 + 4}{80} = -\frac{15}{80} = -\frac{3}{16}
      \]
  \end{enumerate}
\end{soluzione}
\end{esercizio}

\condnewpage

\begin{esercizio}{Calcoli con parentesi}
Esegui i seguenti calcoli progressivi con le parentesi, semplificando le frazioni e riducendo i risultati ai minimi termini.

\vspace{0.15cm}
\begin{tcolorbox}[
  colback=blue!3,
  colframe=mainblue,
  title={Ordine delle parentesi},
  fonttitle=\bfseries\small,
  arc=1.5mm,
  boxrule=0.8pt,
  left=2mm, right=2mm, top=1.5mm, bottom=1.5mm
]
\centering\small
Prima risolvi le \textbf{Tonde} $\mathbf{(\quad)}$, poi le \textbf{Quadre} $\mathbf{[\quad]}$, infine le \textbf{Graffe} $\mathbf{\{\quad\}}$.
\end{tcolorbox}
\vspace{0.2cm}

\renewcommand{\arraystretch}{2.5}
\begin{tabularx}{\linewidth}{XX}
  a) $\frac{21}{24} - \left(\frac{15}{36} - 1 + \frac{7}{12}\right) = \fillbox{2.2cm}$ & b) $\frac{48}{36} - \left(\frac{55}{44} - \frac{6}{15}\right) = \fillbox{2.2cm}$ \\
  c) $8 - \left(\frac{85}{15} + \frac{5}{6}\right) = \fillbox{2.2cm}$ & d) $\frac{50}{40} - \left(\frac{12}{18} + \frac{8}{32}\right) = \fillbox{2.2cm}$ \\
  e) $\frac{45}{72} - \left(\frac{32}{48} - \frac{10}{25}\right) - \frac{1}{40} = \fillbox{2.2cm}$ & f) $\frac{17}{12} + \frac{55}{66} - \left(2 - \frac{12}{16}\right) = \fillbox{2.2cm}$ \\
  g) $\left(\frac{4}{5} - \frac{4}{12}\right) - \frac{1}{3} + \frac{7}{15} = \fillbox{2.2cm}$ & \\
\end{tabularx}
\vspace{0.2cm}

\responsespace{3.8cm}

\begin{soluzione}
  Svolgimento passaggi:
  \begin{enumerate}[label=\alph*)]
    \item Semplificazione e parentesi: Semplifichiamo $\frac{21}{24} = \frac{7}{8}$ e $\frac{15}{36} = \frac{5}{12}$:
      \[
      \frac{7}{8} - \left(\frac{5}{12} - 1 + \frac{7}{12}\right) = \frac{7}{8} - \left(\frac{5+7}{12} - 1\right) = \frac{7}{8} - \left(\frac{12}{12} - 1\right) = \frac{7}{8} - (1-1) = \frac{7}{8}
      \]
    \item Semplificazione e parentesi: Semplifichiamo $\frac{48}{36} = \frac{4}{3}$, $\frac{55}{44} = \frac{5}{4}$ e $\frac{6}{15} = \frac{2}{5}$:
      \[
      \frac{4}{3} - \left(\frac{5}{4} - \frac{2}{5}\right) = \frac{4}{3} - \left(\frac{25-8}{20}\right) = \frac{4}{3} - \frac{17}{20} = \frac{80-51}{60} = \frac{29}{60}
      \]
    \item Semplificazione e parentesi: Semplifichiamo $\frac{85}{15} = \frac{17}{3}$:
      \[
      8 - \left(\frac{17}{3} + \frac{5}{6}\right) = 8 - \left(\frac{34+5}{6}\right) = 8 - \frac{39}{6} = 8 - \frac{13}{2} = \frac{16-13}{2} = \frac{3}{2}
      \]
    \item Semplificazione e parentesi: Semplifichiamo $\frac{50}{40} = \frac{5}{4}$, $\frac{12}{18} = \frac{2}{3}$ e $\frac{8}{32} = \frac{1}{4}$:
      \[
      \frac{5}{4} - \left(\frac{2}{3} + \frac{1}{4}\right) = \frac{5}{4} - \left(\frac{8+3}{12}\right) = \frac{5}{4} - \frac{11}{12} = \frac{15-11}{12} = \frac{4}{12} = \frac{1}{3}
      \]
    \item Semplificazione e mcm (120): Semplifichiamo $\frac{45}{72} = \frac{5}{8}$, $\frac{32}{48} = \frac{2}{3}$ e $\frac{10}{25} = \frac{2}{5}$:
      \[
      \frac{5}{8} - \left(\frac{2}{3} - \frac{2}{5}\right) - \frac{1}{40} = \frac{5}{8} - \frac{4}{15} - \frac{1}{40} = \frac{5 \cdot 15 - 4 \cdot 8 - 1 \cdot 3}{120} = \frac{75 - 32 - 3}{120} = \frac{40}{120} = \frac{1}{3}
      \]
    \item Semplificazione e parentesi: Semplifichiamo $\frac{55}{66} = \frac{5}{6}$ e $\frac{12}{16} = \frac{3}{4}$:
      \[
      \frac{17}{12} + \frac{5}{6} - \left(2 - \frac{3}{4}\right) = \frac{17+10}{12} - \frac{5}{4} = \frac{27}{12} - \frac{15}{12} = \frac{12}{12} = 1
      \]
    \item Semplificazione e parentesi: Semplifichiamo $\frac{4}{12} = \frac{1}{3}$:
      \[
      \left(\frac{4}{5} - \frac{1}{3}\right) - \frac{1}{3} + \frac{7}{15} = \frac{12-5}{15} - \frac{5}{15} + \frac{7}{15} = \frac{7-5+7}{15} = \frac{9}{15} = \frac{3}{5}
      \]
  \end{enumerate}
\end{soluzione}
\end{esercizio}

\condnewpage

\section*{3. Espressioni con frazioni}

\begin{esercizio}{Espressioni con frazioni}
Risolvi le seguenti espressioni con frazioni, rispettando l'ordine delle parentesi. Ricordati di semplificare il risultato finale.

\begin{enumerate}[label=\alph*)]
  \item $\frac{2}{3} + \frac{5}{4} - \left(\frac{3}{6} + 3 - \frac{7}{4}\right) + \left[\frac{5}{3} - \left(\frac{1}{6} + \frac{3}{4} + 1 - \frac{10}{12}\right)\right]$
  \responsespace{3.0cm}

  \item $\left[2 - \frac{2}{3} + \frac{5}{12} - \left(\frac{5}{6} + \frac{7}{12} - \frac{5}{12}\right)\right] + \left(\frac{1}{2} + 6 + \frac{3}{12}\right) - 2$
  \responsespace{3.0cm}

  \item $\frac{4}{3} + \frac{7}{9} - \left(\frac{2}{9} + 1 + \frac{3}{15} + \frac{6}{45}\right) - \left[\frac{18}{15} - \left(1 + \frac{1}{9} - \frac{1}{15}\right)\right]$
  \responsespace{3.0cm}

  \item $\frac{10}{3} + \frac{3}{4} - \left(\frac{7}{12} - \frac{1}{2} + 1\right) - \left[\frac{7}{12} + 3 + \frac{1}{2} - \left(\frac{8}{3} - \frac{9}{4} + \frac{4}{6}\right)\right]$
  \responsespace{3.0cm}

  \item $\frac{11}{12} + \left[ \left( \frac{11}{4} - \frac{7}{5} \right) - \left( 1 - \frac{1}{3} - \frac{1}{2} \right) - \left( \frac{8}{5} - \frac{7}{6} \right) \right]$
  \responsespace{3.0cm}

  \item $4 - \left[ \frac{9}{12} + \left( \frac{5}{6} - \frac{2}{3} \right) - \left( \frac{7}{5} - \frac{4}{3} \right) \right] - \left( 1 - \frac{2}{5} \right)$
  \responsespace{3.0cm}

  \item $\left\{ \left[ \left( \frac{1}{4} + \frac{5}{6} + \frac{13}{2} \right) - \frac{79}{12} \right] + \left[ \left( \frac{11}{4} + \frac{3}{8} + \frac{7}{6} \right) - \frac{101}{24} \right] \right\} - \frac{1}{12}$
  \responsespace{3.0cm}
\end{enumerate}

\begin{soluzione}
  Svolgimento passaggi:
  \begin{enumerate}[label=\alph*)]
    \item Semplifichiamo $\frac{3}{6} = \frac{1}{2}$ nella prima tonda e la calcoliamo:
      \[
      \frac{1}{2} + 3 - \frac{7}{4} = \frac{2 + 12 - 7}{4} = \frac{7}{4}
      \]
      Semplifichiamo $\frac{10}{12} = \frac{5}{6}$ nella tonda interna alla quadra e la calcoliamo:
      \[
      \frac{1}{6} + \frac{3}{4} + 1 - \frac{5}{6} = \frac{3}{4} + 1 - \frac{2}{3} = \frac{9 + 12 - 8}{12} = \frac{13}{12}
      \]
      Risolviamo la parentesi quadra:
      \[
      \left[ \frac{5}{3} - \frac{13}{12} \right] = \frac{20-13}{12} = \frac{7}{12}
      \]
      Eseguiamo il calcolo finale:
      \[
      \frac{2}{3} + \frac{5}{4} - \frac{7}{4} + \frac{7}{12} = \frac{2}{3} - \frac{1}{2} + \frac{7}{12} = \frac{8-6+7}{12} = \frac{9}{12} = \frac{3}{4}
      \]
    \item Svolgiamo la prima tonda interna alla quadra:
      \[
      \frac{5}{6} + \frac{7}{12} - \frac{5}{12} = \frac{5}{6} + \frac{2}{12} = \frac{5}{6} + \frac{1}{6} = \frac{6}{6} = 1
      \]
      Sostituiamo e risolviamo la parentesi quadra:
      \[
      \left[ 2 - \frac{2}{3} + \frac{5}{12} - 1 \right] = \left[ 1 - \frac{2}{3} + \frac{5}{12} \right] = \left[ \frac{1}{3} + \frac{5}{12} \right] = \frac{4+5}{12} = \frac{9}{12} = \frac{3}{4}
      \]
      Eseguiamo il calcolo finale:
      \[
      \frac{3}{4} + \frac{27}{4} - 2 = \frac{30}{4} - 2 = \frac{15}{2} - 2 = \frac{15-4}{2} = \frac{11}{2}
      \]
    \item Semplifichiamo le frazioni nella prima tonda: $\frac{3}{15} = \frac{1}{5}$ e $\frac{6}{45} = \frac{2}{15}$.
      \[
      \frac{2}{9} + 1 + \frac{1}{5} + \frac{2}{15} = \frac{10 + 45 + 9 + 6}{45} = \frac{70}{45} = \frac{14}{9}
      \]
      Risolviamo la tonda interna alla quadra:
      \[
      1 + \frac{1}{9} - \frac{1}{15} = \frac{45 + 5 - 3}{45} = \frac{47}{45}
      \]
      Eseguiamo il calcolo finale:
      \[
      \frac{4}{3} + \frac{7}{9} - \frac{14}{9} - \frac{7}{45} = \frac{4}{3} - \frac{7}{9} - \frac{7}{45} = \frac{60 - 35 - 7}{45} = \frac{18}{45} = \frac{2}{5}
      \]
    \item Risolviamo le due parentesi tonde:
      \[
      \frac{7}{12} - \frac{1}{2} + 1 = \frac{7-6+12}{12} = \frac{13}{12}
      \]
      Semplificando $\frac{4}{6} = \frac{2}{3}$, risolviamo la seconda tonda:
      \[
      \frac{8}{3} - \frac{9}{4} + \frac{2}{3} = \frac{10}{3} - \frac{9}{4} = \frac{40-27}{12} = \frac{13}{12}
      \]
      Sostituiamo nella parentesi quadra:
      \[
      \left[ \frac{7}{12} + 3 + \frac{1}{2} - \frac{13}{12} \right] = \left[ \frac{7 - 13}{12} + 3 + \frac{1}{2} \right] = \left[ -\frac{6}{12} + 3 + \frac{1}{2} \right] = \left[ -\frac{1}{2} + 3 + \frac{1}{2} \right] = 3
      \]
      Eseguiamo il calcolo finale:
      \[
      \frac{10}{3} + \frac{3}{4} - \frac{13}{12} - 3 = \frac{40 + 9 - 13 - 36}{12} = \frac{49-49}{12} = 0
      \]
    \item Svolgiamo prima le tre parentesi tonde:
      1. $\frac{11}{4} - \frac{7}{5} = \frac{55-28}{20} = \frac{27}{20}$ \\
      2. $1 - \frac{1}{3} - \frac{1}{2} = \frac{6-2-3}{6} = \frac{1}{6}$ \\
      3. $\frac{8}{5} - \frac{7}{6} = \frac{48-35}{30} = \frac{13}{30}$ \\
      Inseriamo i risultati nella quadra e troviamo il denominatore comune (60):
      \[
      \frac{11}{12} + \left[ \frac{27}{20} - \frac{1}{6} - \frac{13}{30} \right] = \frac{11}{12} + \left[ \frac{81-10-26}{60} \right] = \frac{11}{12} + \frac{45}{60} = \frac{11}{12} + \frac{3}{4} = \frac{11+9}{12} = \frac{20}{12} = \frac{5}{3}
      \]
    \item Svolgiamo prima le tre parentesi tonde:
      1. $\frac{5}{6} - \frac{2}{3} = \frac{5-4}{6} = \frac{1}{6}$ \\
      2. $\frac{7}{5} - \frac{4}{3} = \frac{21-20}{15} = \frac{1}{15}$ \\
      3. $1 - \frac{2}{5} = \frac{3}{5}$ \\
      Inseriamo i risultati nella quadra e troviamo il denominatore comune (60):
      \[
      4 - \left[ \frac{9}{12} + \frac{1}{6} - \frac{1}{15} \right] - \frac{3}{5} = 4 - \left[ \frac{45+10-4}{60} \right] - \frac{3}{5} = 4 - \frac{51}{60} - \frac{3}{5} = 4 - \frac{17}{20} - \frac{3}{5} = \frac{80-17-12}{20} = \frac{51}{20}
      \]
    \item Risolviamo le due parentesi tonde interne:
      1. $\frac{1}{4} + \frac{5}{6} + \frac{13}{2} = \frac{3+10+78}{12} = \frac{91}{12}$ \\
      2. $\frac{11}{4} + \frac{3}{8} + \frac{7}{6} = \frac{66+9+28}{24} = \frac{103}{24}$ \\
      Risolviamo le due quadre:
      * Prima quadra: $\left[ \frac{91}{12} - \frac{79}{12} \right] = \frac{12}{12} = 1$ \\
      * Seconda quadra: $\left[ \frac{103}{24} - \frac{101}{24} \right] = \frac{2}{24} = \frac{1}{12}$ \\
      Risolviamo la graffa ed eseguiamo l'operazione finale:
      \[
      \left\{ 1 + \frac{1}{12} \right\} - \frac{1}{12} = \frac{13}{12} - \frac{1}{12} = \frac{12}{12} = 1
      \]
  \end{enumerate}
\end{soluzione}
\end{esercizio}

\condnewpage

\begin{esercizio}{La torta di compleanno di Giulia}
Per festeggiare il suo compleanno con gli amici, Giulia prepara una grande torta al cioccolato che pesa $1{,}2\unita{kg}$. Durante la festa:
\begin{itemize}
  \item le ragazze presenti mangiano complessivamente i $\frac{2}{5}$ della torta;
  \item i ragazzi presenti mangiano complessivamente il $\frac{1}{3}$ della torta;
  \item l'insegnante di matematica, invitata alla festa, ne mangia $\frac{1}{10}$.
\end{itemize}

\domanda{a) Quale frazione di torta è stata mangiata in totale?}\\
\domanda{b) Quale frazione di torta rimane per i genitori di Giulia?}\\
\domanda{c) Quanti grammi di torta sono rimasti per i genitori?}

\responsespace{5.5cm}

\begin{soluzione}
  \textbf{Rappresentazione visiva con i blocchetti (totale diviso in 30 parti):}
  \begin{center}
  \begin{tikzpicture}[scale=1.2, every node/.style={scale=0.85}]
    \fill[blue!15!white] (0,0) rectangle (3.0, 0.6); % 12 blocks for Girls (2/5 = 12/30)
    \fill[orange!20!white] (3.0,0) rectangle (5.5, 0.6); % 10 blocks for Boys (1/3 = 10/30)
    \fill[red!15!white] (5.5,0) rectangle (6.25, 0.6); % 3 blocks for Teacher (1/10 = 3/30)
    \fill[green!15!white] (6.25,0) rectangle (7.5, 0.6); % 5 blocks for Parents (1/6 = 5/30)
    \draw[thick, mainblue] (0,0) rectangle (7.5, 0.6);
    \foreach \x in {1,...,29} {
      \draw[linegray, thin] (0.25*\x, 0) -- (0.25*\x, 0.6);
    }
    \node at (1.5, 0.3) {Ragazze ($\frac{12}{30}$)};
    \node at (4.25, 0.3) {Ragazzi ($\frac{10}{30}$)};
    \draw[<-] (5.875, 0.65) -- (5.875, 1.1) node[above] {Prof ($\frac{3}{30}$)};
    \draw[<-] (6.875, -0.05) -- (6.875, -0.55) node[below] {Genitori ($\frac{5}{30}$)};
  \end{tikzpicture}
  \end{center}

  \textbf{Passaggi risolutivi:}
  \begin{enumerate}[label=\alph*)]
    \item \textbf{Frazione totale consumata:} Sommiamo le frazioni mangiate dagli invitati. Il comune denominatore tra $5$, $3$ e $10$ è $30$:
      \[
      \frac{2}{5} + \frac{1}{3} + \frac{1}{10} = \frac{12}{30} + \frac{10}{30} + \frac{3}{30} = \frac{12+10+3}{30} = \frac{25}{30} = \frac{5}{6}
      \]
      In totale è stata consumata la frazione $\frac{5}{6}$ della torta.
    \item \textbf{Frazione rimasta:} Sottraiamo la frazione consumata dall'intera torta ($1$):
      \[
      1 - \frac{5}{6} = \frac{6}{6} - \frac{5}{6} = \frac{1}{6}
      \]
      Ai genitori rimane la frazione $\frac{1}{6}$ della torta.
    \item \textbf{Quantità in grammi:}
      \begin{itemize}
        \item Convertiamo il peso della torta in grammi: $1{,}2\unita{kg} = 1200\unita{g}$.
        \item Calcoliamo la parte rimasta:
          \[
          1200\unita{g} \cdot \frac{1}{6} = \frac{1200}{6}\unita{g} = 200\unita{g}
          \]
      \end{itemize}
  \end{enumerate}
\end{soluzione}
\end{esercizio}

\condnewpage


\begin{esercizio}{La spesa al mercato}
La mamma va al mercato per fare la spesa settimanale con una somma iniziale di denaro. Spende:
\begin{itemize}
  \item i $\frac{2}{7}$ del denaro totale dal fruttivendolo per acquistare frutta e verdura;
  \item i $\frac{3}{14}$ del denaro totale dal panettiere per pane e focaccia.
\end{itemize}

\domanda{a) Quale frazione del denaro iniziale ha speso complessivamente la mamma?}\\
\domanda{b) Quale frazione di denaro le rimane nel portafoglio?}\\
\domanda{c) Se inizialmente la mamma aveva nel portafoglio $56$ euro, quanti euro le rimangono alla fine della spesa?}

\responsespace{5.5cm}

\begin{soluzione}
  \textbf{Rappresentazione visiva con i blocchetti (totale diviso in 14 parti):}
  \begin{center}
  \begin{tikzpicture}[scale=1.3, every node/.style={scale=0.85}]
    \fill[blue!15!white] (0,0) rectangle (2.0, 0.6); % 4 blocks for Fruttivendolo (2/7 = 4/14)
    \fill[orange!20!white] (2.0,0) rectangle (3.5, 0.6); % 3 blocks for Panettiere (3/14)
    \fill[green!15!white] (3.5,0) rectangle (7.0, 0.6); % 7 blocks Rimanenti (1/2 = 7/14)
    \draw[thick, mainblue] (0,0) rectangle (7.0, 0.6);
    \foreach \x in {1,...,13} {
      \draw[linegray, thin] (0.5*\x, 0) -- (0.5*\x, 0.6);
    }
    \draw[<-] (1.0, -0.05) -- (1.0, -0.55) node[below] {Fruttivendolo ($\frac{4}{14}$)};
    \draw[<-] (2.75, 0.65) -- (2.75, 1.1) node[above] {Panettiere ($\frac{3}{14}$)};
    \node at (5.25, 0.3) {Rimanente ($\frac{1}{2}$)};
  \end{tikzpicture}
  \end{center}

  \textbf{Passaggi:}
  \begin{enumerate}[label=\alph*)]
    \item \textbf{Frazione di denaro spesa:} Sommiamo le due spese parziali. Il comune denominatore tra $7$ e $14$ è $14$:
      \[
      \frac{2}{7} + \frac{3}{14} = \frac{4}{14} + \frac{3}{14} = \frac{4+3}{14} = \frac{7}{14} = \frac{1}{2}
      \]
      Complessivamente la mamma ha speso la frazione $\frac{1}{2}$ (esattamente la metà) del suo budget.
    \item \textbf{Frazione di denaro rimanente:} Sottraiamo la frazione spesa dal totale ($1$):
      \[
      1 - \frac{1}{2} = \frac{2}{2} - \frac{1}{2} = \frac{1}{2}
      \]
      Le rimane la frazione $\frac{1}{2}$ della somma iniziale.
    \item \textbf{Calcolo dell'importo rimanente in euro (budget = $56$ euro):}
      \begin{itemize}
        \item Calcoliamo la metà di $56$ euro:
          \[
          56\text{ euro} \cdot \frac{1}{2} = \frac{56}{2}\text{ euro} = 28\text{ euro}
          \]
      \end{itemize}
  \end{enumerate}
\end{soluzione}
\end{esercizio}

\condnewpage


\begin{esercizio}{La gita scolastica}
Durante la gita scolastica di primavera della classe 2B, il viaggio totale per raggiungere il parco naturale è suddiviso nel seguente modo:
\begin{itemize}
  \item i $\frac{3}{5}$ del tragitto totale vengono percorsi in treno;
  \item il $\frac{1}{4}$ del tragitto totale viene percorso in autobus;
  \item la parte rimanente del viaggio viene percorsa a piedi lungo un sentiero naturalistico.
\end{itemize}

\domanda{a) Quale frazione del viaggio totale è coperta da treno e autobus insieme?}\\
\domanda{b) Quale frazione del viaggio viene percorsa a piedi?}\\
\domanda{c) Se il tragitto complessivo è lungo $80\unita{km}$, quanti chilometri percorrono a piedi?}

\responsespace{5.5cm}

\begin{soluzione}
  \textbf{Rappresentazione visiva con i blocchetti (totale diviso in 20 parti):}
  \begin{center}
  \begin{tikzpicture}[scale=1.2, every node/.style={scale=0.85}]
    \fill[blue!15!white] (0,0) rectangle (4.8, 0.6); % 12 blocks for Treno (3/5 = 12/20)
    \fill[orange!20!white] (4.8,0) rectangle (6.8, 0.6); % 5 blocks for Autobus (1/4 = 5/20)
    \fill[green!15!white] (6.8,0) rectangle (8.0, 0.6); % 3 blocks a piedi (3/20)
    \draw[thick, mainblue] (0,0) rectangle (8.0, 0.6);
    \foreach \x in {1,...,19} {
      \draw[linegray, thin] (0.4*\x, 0) -- (0.4*\x, 0.6);
    }
    \node at (2.4, 0.3) {Treno ($\frac{12}{20}$)};
    \node at (5.8, 0.3) {Autobus ($\frac{5}{20}$)};
    \draw[<-] (7.4, -0.05) -- (7.4, -0.55) node[below] {A piedi ($\frac{3}{20}$)};
  \end{tikzpicture}
  \end{center}

  \textbf{Passaggi:}
  \begin{enumerate}[label=\alph*)]
    \item \textbf{Frazione treno + autobus:} Sommiamo le due frazioni del viaggio. Il comune denominatore tra $5$ e $4$ è $20$:
      \[
      \frac{3}{5} + \frac{1}{4} = \frac{12}{20} + \frac{5}{20} = \frac{17}{20}
      \]
      Treno e autobus coprono insieme i $\frac{17}{20}$ del percorso totale.
    \item \textbf{Frazione a piedi:} Sottraiamo dal totale della gita ($1$):
      \[
      1 - \frac{17}{20} = \frac{20}{20} - \frac{17}{20} = \frac{3}{20}
      \]
      La parte a piedi corrisponde ai $\frac{3}{20}$ del viaggio complessivo.
    \item \textbf{Distanza a piedi:}
      \begin{itemize}
        \item Calcoliamo la distanza in km sapendo che l'intero tragitto misura $80\unita{km}$:
          \[
          80\unita{km} \cdot \frac{3}{20} = \frac{80 \cdot 3}{20}\unita{km} = 4 \cdot 3\unita{km} = 12\unita{km}
          \]
      \end{itemize}
  \end{enumerate}
\end{soluzione}
\end{esercizio}

\condnewpage


\section*{5. Sfide}

\begin{esercizio}{Come arriviamo a scuola?}
In una classe di 24 alunni, si analizzano i mezzi che gli studenti usano per raggiungere la scuola:
\begin{itemize}
  \item la metà degli alunni ($\frac{1}{2}$) raggiunge la scuola a piedi;
  \item un terzo degli alunni ($\frac{1}{3}$) usa di solito la bicicletta;
  \item i restanti alunni utilizzano i mezzi pubblici.
\end{itemize}

\domanda{a) Quale frazione degli alunni rappresenta il terzo gruppo che usa i mezzi pubblici?}\\
\domanda{b) Calcola quanti alunni fanno parte di ciascuno dei tre gruppi.}

\responsespace{5.5cm}

\begin{soluzione}
  \textbf{Rappresentazione visiva con i blocchetti (totale diviso in 6 parti):}
  \begin{center}
  \begin{tikzpicture}[scale=1.4, every node/.style={scale=0.85}]
    \fill[blue!15!white] (0,0) rectangle (3.0, 0.6); % 3 blocks for A piedi (1/2 = 3/6)
    \fill[orange!20!white] (3.0,0) rectangle (5.0, 0.6); % 2 blocks for Bicicletta (1/3 = 2/6)
    \fill[green!15!white] (5.0,0) rectangle (6.0, 0.6); % 1 block for Mezzi pubblici (1/6)
    \draw[thick, mainblue] (0,0) rectangle (6.0, 0.6);
    \foreach \x in {1,...,5} {
      \draw[linegray, thin] (1.0*\x, 0) -- (1.0*\x, 0.6);
    }
    \node at (1.5, 0.3) {A piedi ($\frac{3}{6}$)};
    \node at (4.0, 0.3) {Bici ($\frac{2}{6}$)};
    \draw[<-] (5.5, -0.05) -- (5.5, -0.55) node[below] {Mezzi ($\frac{1}{6}$)};
  \end{tikzpicture}
  \end{center}

  \textbf{Passaggi risolutivi:}
  \begin{enumerate}[label=\alph*)]
    \item \textbf{Frazione dei mezzi pubblici:} Sommiamo le frazioni di chi va a piedi e in bici. Il comune denominatore tra $2$ e $3$ è $6$:
      \[
      \frac{1}{2} + \frac{1}{3} = \frac{3}{6} + \frac{2}{6} = \frac{5}{6}
      \]
      La parte restante che usa i mezzi pubblici è:
      \[
      1 - \frac{5}{6} = \frac{6}{6} - \frac{5}{6} = \frac{1}{6}
      \]
      Il terzo gruppo corrisponde alla frazione $\frac{1}{6}$ dell'intera classe.
    \item \textbf{Numero di alunni per gruppo (totale = 24 allievi):}
      \begin{itemize}
        \item \textbf{A piedi:} $24 \cdot \frac{1}{2} = 12$ alunni.
        \item \textbf{In bicicletta:} $24 \cdot \frac{1}{3} = 8$ alunni.
        \item \textbf{Mezzi pubblici:} $24 \cdot \frac{1}{6} = 4$ alunni.
      \end{itemize}
      Verifica: $12 + 8 + 4 = 24$ alunni.
  \end{enumerate}
\end{soluzione}
\end{esercizio}

\condnewpage

\begin{esercizio}{La staffetta di atletica}
Tre studenti Matteo, Sofia e Luca partecipano ad una corsa a staffetta sulla pista della scuola, che è lunga in totale $400\unita{m}$. La ripartizione della corsa è la seguente:
\begin{itemize}
  \item Sofia corre una distanza pari al \textbf{doppio} di quella corsa da Matteo;
  \item Luca corre la parte rimanente della pista, che corrisponde esattamente ad un quarto ($\frac{1}{4}$) del totale.
\end{itemize}

\domanda{a) Osserva lo schema a blocchetti qui sotto: spiega a parole perché la corsa di Sofia corrisponde a metà dell'intero percorso.}\\
\domanda{b) Calcola quanti metri corre ciascuno dei tre studenti.}

\vspace{0.4cm}
\begin{center}
\begin{tikzpicture}[scale=1.2, every node/.style={scale=0.85}]
  % Overall block representation
  \draw[thick, mainblue] (0,0) rectangle (8, 0.6);
  
  % Internal divisions
  \draw[thick, mainblue] (2,0) -- (2, 0.6);
  \draw[thick, mainblue] (6,0) -- (6, 0.6);
  \draw[dashed, linegray] (4,0) -- (4, 0.6); % Sofia's half line for reference
  
  % Annotations below the bar for student/solution
  \ifsoluzioni
    \fill[blue!15!white] (0,0) rectangle (2, 0.6);
    \fill[orange!20!white] (2,0) rectangle (6, 0.6);
    \fill[green!15!white] (6,0) rectangle (8, 0.6);
    
    \draw[thick, mainblue] (0,0) rectangle (8, 0.6);
    \draw[thick, mainblue] (2,0) -- (2, 0.6);
    \draw[thick, mainblue] (6,0) -- (6, 0.6);
    
    \node at (1, 0.3) {Matteo};
    \node at (1,-0.35) {$\frac{1}{4}$ ($100\unita{m}$)};
    
    \node at (4, 0.3) {Sofia};
    \node at (4,-0.35) {$\frac{1}{2}$ ($200\unita{m}$)};
    
    \node at (7, 0.3) {Luca};
    \node at (7,-0.35) {$\frac{1}{4}$ ($100\unita{m}$)};
  \else
    \node at (1, 0.3) {Matteo};
    \node at (4, 0.3) {Sofia};
    \node at (7, 0.3) {Luca};
  \fi
\end{tikzpicture}
\end{center}

\responsespace{5.5cm}

\begin{soluzione}
  \textbf{Svolgimento e ragionamento:}
  \begin{enumerate}[label=\alph*)]
    \item \textbf{Ragionamento con i blocchetti (Spiegazione):} Poiché a Luca spetta $\frac{1}{4}$ del tracciato (l'ultimo blocco), la parte restante per Matteo e Sofia insieme è $1 - \frac{1}{4} = \frac{3}{4}$ della pista.
      Sapendo che Sofia corre il doppio di Matteo, possiamo dividere questa parte restante ($\frac{3}{4}$) in $3$ blocchi di uguale lunghezza: $1$ per Matteo e $2$ per Sofia.
      \begin{itemize}
        \item Ciascun blocco corrisponde a: $\frac{3}{4} : 3 = \frac{1}{4}$ della pista totale.
        \item Di conseguenza, Matteo corre $\frac{1}{4}$, Sofia corre $2$ blocchi cioè $2 \cdot \frac{1}{4} = \frac{2}{4} = \frac{1}{2}$ (che corrisponde esattamente a metà dell'intero percorso!), e Luca corre $\frac{1}{4}$.
      \end{itemize}
    \item \textbf{Calcolo delle distanze reali (totale = $400\unita{m}$):}
      Utilizzando lo schema a blocchi o le frazioni ricavate:
      \begin{itemize}
        \item \textbf{Matteo:} corre $\frac{1}{4}$ del percorso: $400\unita{m} \cdot \frac{1}{4} = 100\unita{m}$.
        \item \textbf{Sofia:} corre $\frac{1}{2}$ del percorso: $400\unita{m} \cdot \frac{1}{2} = 200\unita{m}$.
        \item \textbf{Luca:} corre $\frac{1}{4}$ del percorso: $400\unita{m} \cdot \frac{1}{4} = 100\unita{m}$.
      \end{itemize}
      Verifica: $100\unita{m} + 200\unita{m} + 100\unita{m} = 400\unita{m}$.
  \end{enumerate}
\end{soluzione}
\end{esercizio}

\condnewpage

\begin{esercizio}{Il livello del serbatoio d'acqua}
Un serbatoio per la raccolta dell'acqua piovana dell'orto scolastico è riempito fino a una frazione ignota $F$ della sua capacità totale. Durante una notte di pioggia avvengono due fatti:
\begin{itemize}
  \item l'acqua piovana fa aumentare l'acqua già presente nel serbatoio di \textbf{un quarto del suo stesso volume} (quindi si aggiunge $\frac{1}{4}$ dell'acqua già presente);
  \item il canale di gronda del tetto convoglia nel serbatoio una quantità d'acqua pari a \textbf{un decimo ($\frac{1}{10}$) della capacità totale} del serbatoio.
\end{itemize}
Al mattino, il serbatoio si trova riempito esattamente per i tre quinti ($\frac{3}{5}$) della sua capacità totale.

\domanda{a) Trova la frazione iniziale $F$ d'acqua nel serbatoio prima della pioggia (puoi aiutarti impostando un'uguaglianza descrittiva con le frazioni).}\\
\domanda{b) Se il serbatoio contiene in totale $2000$ litri quando è pieno, quanti litri d'acqua c'erano all'inizio? Quanti litri sono stati aggiunti complessivamente durante la notte?}

\responsespace{6.5cm}

\begin{soluzione}
  \textbf{Svolgimento e ragionamento:}
  \begin{enumerate}[label=\alph*)]
    \item \textbf{Calcolo della frazione iniziale F:}
      Impostiamo una semplice uguaglianza a parole e poi con le frazioni:
      \begin{itemize}
        \item Frazione iniziale d'acqua: $F$.
        \item Aumento dovuto alla pioggia diretta: $\frac{1}{4}$ del volume preesistente, ossia $\frac{1}{4} \cdot F$. L'acqua iniziale più l'aumento diretto diventa:
          \[
          F + \frac{1}{4}F = \frac{5}{4}F
          \]
        \item Aggiunta costante della grondaia: $\frac{1}{10}$ della capacità totale.
        \item Frazione finale al mattino: $\frac{3}{5}$.
      \end{itemize}
      L'uguaglianza che descrive la situazione è:
      \[
      \frac{5}{4}F + \frac{1}{10} = \frac{3}{5}
      \]
      Isoliamo il termine $F$ portando $\frac{1}{10}$ a destra. Il comune denominatore tra $5$ e $10$ è $10$:
      \[
      \frac{5}{4}F = \frac{3}{5} - \frac{1}{10} \implies \frac{5}{4}F = \frac{6}{10} - \frac{1}{10} \implies \frac{5}{4}F = \frac{5}{10} = \frac{1}{2}
      \]
      Ora troviamo $F$ moltiplicando la frazione per il reciproco $\frac{4}{5}$:
      \[
      F = \frac{1}{2} \cdot \frac{4}{5} = \frac{4}{10} = \frac{2}{5}
      \]
      Prima della pioggia, il serbatoio era pieno per i $\frac{2}{5}$ della sua capacità totale.
    \item \textbf{Calcolo dei litri (capacità totale = $2000\unita{l}$):}
      Utilizzando la frazione iniziale $F = \frac{2}{5}$ e finale $\frac{3}{5}$:
      \begin{itemize}
    \item \textbf{Volume iniziale d'acqua:} $\frac{2}{5} \cdot 2000\unita{l} = 2 \cdot 400\unita{l} = 800\unita{l}$.
        \item \textbf{Volume finale d'acqua:} $\frac{3}{5} \cdot 2000\unita{l} = 3 \cdot 400\unita{l} = 1200\unita{l}$.
        \item \textbf{Litri totali aggiunti:} $1200\unita{l} - 800\unita{l} = 400\unita{l}$.
      \end{itemize}
  \end{enumerate}
\end{soluzione}
\end{esercizio}

\condnewpage

\begin{esercizio}{Il disegno misterioso}
Esegui i seguenti calcoli di riepilogo con le frazioni, riducendo sempre i risultati ai minimi termini. Semplifica dove possibile.

\vspace{0.3cm}
\renewcommand{\arraystretch}{2.25}
\begin{tabularx}{\linewidth}{XX}
  1) $\frac{3}{4} + \frac{2}{5} = \fillbox{2.2cm}$ & 2) $\frac{2}{7} + \frac{1}{3} = \fillbox{2.2cm}$ \\
  3) $\frac{5}{8} - \frac{1}{6} = \fillbox{2.2cm}$ & 4) $\frac{3}{2} - \frac{5}{9} = \fillbox{2.2cm}$ \\
  5) $\frac{4}{5} - \frac{13}{20} = \fillbox{2.2cm}$ & 6) $\frac{1}{2} + \frac{2}{5} - \frac{3}{20} = \fillbox{2.2cm}$ \\
  7) $\frac{5}{6} + \left[ \frac{3}{4} - \frac{1}{3} \right] = \fillbox{2.2cm}$ & 8) $\frac{10}{7} - \left[ \frac{2}{3} - \frac{4}{21} \right] = \fillbox{2.2cm}$ \\
  9) $4 - \left[ \frac{5}{6} + \frac{1}{3} - \frac{3}{4} \right] = \fillbox{2.2cm}$ & 10) $\left[ \frac{7}{6} - \frac{1}{4} \right] - \frac{5}{12} = \fillbox{2.2cm}$ \\
  11) $\left[ \frac{10}{16} + \frac{8}{20} \right] - \left[ \frac{6}{12} - \frac{5}{25} \right] = \fillbox{2.2cm}$ & 12) $\frac{10}{50} + \frac{6}{36} + \frac{8}{48} = \fillbox{2.2cm}$ \\
  13) $\frac{24}{3} + \frac{9}{27} = \fillbox{2.2cm}$ & 14) $4 - \left[ \frac{4}{5} + \frac{7}{10} + \frac{5}{6} \right] = \fillbox{2.2cm}$ \\
\end{tabularx}

\condnewpage

\domanda{La vignetta misteriosa}\\
Colorando nella figura sottostante solo i settori che contengono i risultati esatti dei 14 calcoli della pagina precedente, scoprirai una vignetta misteriosa!

\vspace{0.5cm}
\immagine[14.0cm]{vignetta}

\responsespace{0.5cm}

\begin{soluzione}
  \textbf{Risultati esatti dei 14 calcoli (ridotti ai minimi termini):}
  \begin{enumerate}[label=\arabic*)]
    \item $\frac{3}{4} + \frac{2}{5} = \frac{15+8}{20} = \frac{23}{20}$
    \item $\frac{2}{7} + \frac{1}{3} = \frac{6+7}{21} = \frac{13}{21}$
    \item $\frac{5}{8} - \frac{1}{6} = \frac{15-4}{24} = \frac{11}{24}$
    \item $\frac{3}{2} - \frac{5}{9} = \frac{27-10}{18} = \frac{17}{18}$
    \item $\frac{4}{5} - \frac{13}{20} = \frac{16-13}{20} = \frac{3}{20}$
    \item $\frac{1}{2} + \frac{2}{5} - \frac{3}{20} = \frac{10+8-3}{20} = \frac{15}{20} = \frac{3}{4}$
    \item $\frac{5}{6} + \left[ \frac{3}{4} - \frac{1}{3} \right] = \frac{5}{6} + \frac{9-4}{12} = \frac{5}{6} + \frac{5}{12} = \frac{10+5}{12} = \frac{15}{12} = \frac{5}{4}$
    \item $\frac{10}{7} - \left[ \frac{2}{3} - \frac{4}{21} \right] = \frac{10}{7} - \frac{14-4}{21} = \frac{10}{7} - \frac{10}{21} = \frac{30-10}{21} = \frac{20}{21}$
    \item $4 - \left[ \frac{5}{6} + \frac{1}{3} - \frac{3}{4} \right] = 4 - \frac{10+4-9}{12} = 4 - \frac{5}{12} = \frac{48-5}{12} = \frac{43}{12}$
    \item $\left[ \frac{7}{6} - \frac{1}{4} \right] - \frac{5}{12} = \frac{14-3}{12} - \frac{5}{12} = \frac{11-5}{12} = \frac{6}{12} = \frac{1}{2}$
    \item $\left[ \frac{10}{16} + \frac{8}{20} \right] - \left[ \frac{6}{12} - \frac{5}{25} \right] = \left[ \frac{5}{8} + \frac{2}{5} \right] - \left[ \frac{1}{2} - \frac{1}{5} \right] = \frac{25+16}{40} - \frac{5-2}{10} = \frac{41}{40} - \frac{3}{10} = \frac{41-12}{40} = \frac{29}{40}$
    \item $\frac{10}{50} + \frac{6}{36} + \frac{8}{48} = \frac{1}{5} + \frac{1}{6} + \frac{1}{6} = \frac{1}{5} + \frac{2}{6} = \frac{1}{5} + \frac{1}{3} = \frac{3+5}{15} = \frac{8}{15}$
    \item $\frac{24}{3} + \frac{9}{27} = 8 + \frac{1}{3} = \frac{24+1}{3} = \frac{25}{3}$
    \item $4 - \left[ \frac{4}{5} + \frac{7}{10} + \frac{5}{6} \right] = 4 - \frac{24+21+25}{30} = \frac{5}{3}$
  \end{enumerate}
\end{soluzione}
\end{esercizio}
